% Related work.

\section{Related work}\label{sec:related}

The paper draws on four strands of prior work: the expressiveness
theory of graph neural networks, the polarisation measurement
literature in network science and computational social science, the
line of graph neural network estimators applied to polarisation and
ideology on social platforms, and the null-model rewiring literature
in network science. This chapter surveys each strand and locates the
paper's contribution.

\subsection{Graph neural network expressiveness}

Message-passing graph neural networks were introduced in a general
form by \citet{GilmerNeuralMessagePassing2017}, unifying earlier
architectures for molecular property prediction. Their expressiveness
was placed against the Weisfeiler-Leman graph isomorphism test by
\citet{MorrisWL2019}: any GNN that follows the message-passing schema
of Definition~\ref{def:mpnn} distinguishes graphs no better than $1$-WL colour
refinement. \citet{XuGIN2019} gave the matching lower bound with the
graph isomorphism network, which uses sum aggregation and multi-layer
perceptron updates to attain $1$-WL expressiveness on countable
feature domains.

The passage to continuous feature domains and to universal
approximation on featured graphs was made by \citet{ChenVillar2019},
who established a two-way equivalence between graph-isomorphism
discrimination and universal approximation of permutation-invariant
continuous functionals on compact featured-graph classes. Their
result is per-graph-size, matching the format of Theorem~\ref{thm:dichotomy}
in this paper. \citet{AzizianLelarge2021} tightened the equivalence
to the form we use: on a compact bounded-size featured-graph space,
the uniform closure of the message-passing GNN class equals exactly
the continuous $1$-WL invariant functionals. Higher-order variants
(hierarchical Weisfeiler-Leman refinements and $k$-GNNs) escape the
$1$-WL bound and gain expressiveness at the cost of computational
scaling with $k$ \citep{MorrisWL2019}.

Architectures that inject subgraph counts as vertex features
\citep{Bouritsas2022} also escape the $1$-WL bound. Subgraph counts
include triangles, so these architectures can compute closure
directly; the impossibility in this paper does not apply to them.
Graph transformers with positional encodings based on the Laplacian
spectrum sit outside the message-passing schema of Definition~\ref{def:mpnn}
entirely.

The graphon-limit theory of GNN expressiveness \citep{KerivenPeyre2020,
RuizChamonRibeiro2021} studies convergence of GNN outputs on
random-graph samples of a graphon and shows universal approximation
of continuous cut-metric functionals in the graphon limit. The
implication for the impossibility of this paper is discussed in
Section~\ref{sec:obstacle}: the ensemble-level entanglement does not
extend to the dense graphon regime, only to the sparse and
structured regimes real interaction graphs inhabit.

\subsection{Polarisation measurement}

The measurement of political polarisation on social platforms is a
mature field. Homophily indices and modularity have long been used to
assess whether interaction graphs align with partisan labels
\citep{Newman2003, Newman2004}. \citet{AdamicGlance2005} applied assortativity-based analyses to
early blog data. \citet{ConoverPolBlogs2011} extended homophily-based
scoring to Twitter interaction graphs.

More structural measures followed. \citet{Garimella18} introduced the
random walk controversy score (Definition~\ref{def:rwc}), which measures the
propensity of random walks starting on one side of a bipartition to
terminate on the same side. \citet{Hohmann2016} and later
\citet{HohmannGED2023} introduced spectral measures culminating in
the generalised Euclidean distance (Definition~\ref{def:ged}), which measures
opinion separation via the Laplacian pseudoinverse and admits a
partition-free formulation.

At the community level, \citet{WallerAnderson2021} released
$150$-dimensional embeddings for ten thousand subreddits, along with
one-dimensional partisan-axis scores per subreddit. Their embedding
procedure is a skip-gram model on user co-comment counts, not a
graph neural network. \citet{Barbera2015} developed a Bayesian
ideal-point model on Twitter follower networks; ideal-point models
lie outside the message-passing class and are unaffected by the
dichotomy of Theorem~\ref{thm:dichotomy}.

None of the polarisation measurement literature before this paper
has, to our knowledge, connected the choice of measure to the
expressiveness class of the estimator applied to it. The taxonomy
of Section~\ref{sec:measures} is the paper's contribution in this
direction.

\subsection{GNN-based polarisation estimation}

A recent line of work applies graph neural networks directly to
polarisation and ideology detection on interaction graphs.

EchoGAE \citep{Alatawi23EchoGAE} is a two-layer graph auto-encoder
whose Echo Chamber Score is a cohesion-versus-separation metric in
the vertex embedding space. The score is unsupervised: it does not
see partisan labels at training. Its authors report positive results
on four Twitter datasets.

InfoVGAE \citep{Li22InfoVGAE} is a variational graph auto-encoder
with disentangled latent axes intended to capture belief dimensions.
The variational objective adds a total-correlation regulariser to
the reconstruction loss; the resulting per-vertex embeddings are
used to score users on inferred belief axes.

TIMME \citep{XiaoTIMME2020} is a relational graph convolutional
network for Twitter ideology detection. It runs over the
heterogeneous graph formed by follow, retweet, mention, reply, and
like edges and outputs a supervised ideology label for each user.

Each of these three estimators is a message-passing graph neural
network in the sense of Definition~\ref{def:mpnn}. All three fall within the
scope of Theorem~\ref{thm:dichotomy}. The empirical audit of
Section~\ref{sec:emp-reddit} instantiates the procedure derived from
Proposition~\ref{prop:reproduce} on EchoGAE; the same procedure applies to
InfoVGAE and TIMME with no change beyond the choice of
per-vertex features.

Beyond these three, several other message-passing approaches to
polarisation and stance detection have been proposed in the recent
literature, with heterogeneous input types and multimodal features.
All fall within the message-passing class as defined here and are
therefore subject to Theorem~\ref{thm:dichotomy}.

\subsection{Null-model rewiring in network science}

Degree-preserving rewiring, sometimes called the configuration model
or the null-model construction, dates to work by \citet{MolloyReed1995}
and \citet{Newman2001} in the study of random graphs with prescribed
degree sequence. The construction has since been used extensively as
a null model for tests of higher-order structure: given a graph
$\gph$, one compares the observed value of a structural quantity
(such as clustering or modularity) against the distribution of the
quantity under uniform rewiring conditional on the degree sequence.
\citet{Karrer2011StochasticBlock} developed the degree-corrected
stochastic block model, extending the null-model framework to
include community structure.

\citet{Karrer2011StochasticBlock} developed the degree-corrected
stochastic block model, which combines degree correction with a
block structure and is standard in null-model tests of community
detection. Block-preserving rewiring is standard in this line;
targeted triangle-destroying rewiring is less common. To our
knowledge, no prior work uses joint degree- and block-preserving
rewiring as an identifiability test for GNN-based estimators. The
rewiring procedure of Section~\ref{sec:emp-reddit} is adapted to that
purpose: it holds the depth-$1$ $1$-WL colour distribution of the
labelled edges exact by construction, and drops higher-order
structure by targeted swaps.

Adversarial perturbations of graph structure for GNN robustness
\citep{Jin2020, Zugner2018} share the technical mechanism (edge
addition, deletion, or swap) with this paper's rewiring but use it
under a different objective: adversarial robustness studies MPNN
sensitivity to worst-case edge perturbations that maximise output
change, whereas the rewiring here is constrained to preserve
$1$-WL statistics that any MPNN can compute at depth $1$.
